\section{Related Work}
\label{sec:related_work}

To contextualize our contributions, this section discusses relevant state-of-the-art methodologies that we will compare against. We focus on advancements in self-supervised learning for time series, the emergence of foundation models for physiological signals, and existing approaches to managing variable input structures, especially concerning topological heterogeneity in the EEG domain and computational efficiency.

\subsection{Self-Supervised Learning Strategies in EEG}
Foundation models for EEG primarily rely on self-supervised learning (SSL) to leverage large unlabeled datasets. Masked signal modeling is a dominant paradigm. BENDR \citep{kostas2021bendr} pioneered this for EEG by adapting masked prediction concepts from speech, applying a contrastive objective to predict masked convolutional features. Subsequent models refined this: BrainBERT \citep{wang2023brainbert} performs masked prediction on channel-independent spectrograms for intracranial electroencephalography (iEEG); EEGFormer \citep{chen2024eegformer} and LaBraM \citep{jiang2024large} predict vector-quantized (VQ) representations of masked patches, learning discrete codebooks; CBraMod \citep{wang2024cbramod} directly reconstructs masked raw signal patches. LUNA employs a similar masked reconstruction objective but applies it after projecting channel information into a unified latent space, requiring the decoder to reconstruct channel-specific details from this compressed representation. 

\subsection{Modeling Spatial Structure and Topology Variation in EEG}
\label{subsec:related_work}
Capturing the spatial relationships between EEG channels is vital but complicated by varying electrode counts and layouts across datasets. Several strategies have been explored in the literature: \newline
\textbf{Channel Independence:} Early approaches and models like BrainBERT \citep{wang2023brainbert} and EEGFormer \citep{chen2024eegformer} process each channel's data independently before potentially combining them later. While inherently handling varying channel numbers, this neglects early modeling of cross-channel interactions. \newline
\textbf{Fixed-Topology Spatial Modeling:} Models like Brant \citep{zhang2023brant} use dedicated spatial encoders alongside temporal ones but assume a consistent channel configuration, limiting cross-dataset generalization. Graph Neural Networks (GNNs) \citep{tang2021self} explicitly model spatial relationships using a predefined adjacency graph, but require mechanisms to handle dynamically changing graph structures when topologies vary. LUNA avoids pre-defined graphs or fixed structures.  \newline
\textbf{Joint Spatio-Temporal Attention:} LaBraM \citep{jiang2024large} flattens channel and patch dimensions into one long sequence, allowing a standard Transformer to learn spatio-temporal dependencies simultaneously. However, this incurs $\bigO((S\cdot C)^2)$ complexity, scaling quadratically with both sequence length/patches (S) and channels (C). CBraMod \citep{wang2024cbramod} and CEReBrO \citep{dimofte2025cerebro} use alternating or parallel spatial and temporal attention mechanisms, reducing complexity to $\bigO(max(S^2, C^2))$ but still scaling quadratically with the dominant dimension. BIOT \citep{yang2023biot} uses linear attention after flattening, improving efficiency but potentially limiting modeling capacity. LUNA differs significantly by performing channel unification first before applying temporal attention with quadratic complexity only on the patch dimension and the much smaller latent dimension Q.  \newline
\textbf{Explicit Topology Mapping:} Some methods explicitly map varying topologies to a canonical representation. MMM \citep{yi2023learning} maps channels to predefined anatomical regions but relies on hand-engineered features (Differential Entropy) rather than raw signals. PopT \citep{chau2024population} aggregates pre-computed channel-independent temporal features using 3D electrode coordinates. While achieving topology invariance, these methods are not fully end-to-end or rely on external information (regions). LUNA learns an end-to-end mapping from raw signals using learned queries without requiring pre-defined structures.
\textbf{Differentiable Channel Reordering.}
Saeed et al.~\citep{saeed2021learning} also use learned attention to reorder/project heterogeneous channels into a fixed space, but under supervised training and without explicit 3D channel encodings. In contrast, LUNA is a self-supervised foundation model trained at scale, integrates explicit spatial (3D) information, and targets topology-agnostic \emph{and} compute-efficient transfer across datasets and tasks.
\subsection{Learned Queries and Efficient Attention for Set Abstraction}
LUNA's core mechanism for topology unification draws inspiration from architectures designed for permutation-invariant processing of set-structured data. Set Transformer \citep{lee2019set} introduced the concept of using a small set of learnable inducing points (queries) and an Induced Set Attention Block to summarize information from a larger input set via cross-attention, reducing the complexity from $\bigO(N^2)$ to $\bigO(M \cdot N)$. PerceiverIO \citep{jaegle2021perceiver} further developed this mechanism, demonstrating its power in creating a fixed-size latent bottleneck capable of handling diverse, variable-sized inputs across different modalities (images, text) and enabling flexible decoding via task-specific output queries.

LUNA adapts this principle specifically for EEG topology invariance. We treat the set of EEG channel features at a given time interval (patch) as the input set. By applying cross-attention between the channel features (as keys/values) and a small number (Q) of learned queries, LUNA projects the variable-channel input onto a fixed-size latent space ($\mathbb{R}^{Q \times E}$). This projection is permutation-invariant with respect to the input channels, thus achieving topology agnosticism. Furthermore, it improves computational efficiency, as the complexity of this step scales linearly with the number of channels. Where MMM relies on predefined regions (and hand-crafted features) and LaBraM/CBraMod rely on quadratic spatial attention after flattening space–time, LUNA first \emph{unifies} variable channel sets with learned queries (with explicit 3D encodings) and \emph{then} applies temporal attention on a fixed-size latent, yielding linear-in-$C$ unification and reduced temporal sequence length. This design specifically targets topology-agnostic scaling and inference efficiency within a self-supervised foundation-model framework, rather than combining prior components unchanged.



